
%%%%%%%%%%%%%%%%%%%%%%%%%%%%%%%%%%%%%%%%%%%%%%%%%%%%%%%%%%%%%
\section{模型的改进与推广}

\subsection{模型的改进}

针对模型缺点，可从以下方向改进：引入多周期动态规划或贝叶斯更新，利用每日新观测修正后续需求预测，改善单周期独立性假设；如获得库存和缺货数据，将建模目标从"销售量"修正为"需求量"；如获得天气和促销数据，作为Prophet的外生回归项改善预测精度；对Q3引入单品间的交叉弹性，考虑替代和互补效应；对稀疏单品采用分层贝叶斯模型，借用同品类高销量单品的信息改善估计。

\subsection{模型的推广}

本文模型框架适用于其他易腐生鲜品（水果、乳制品、肉类、烘焙食品）的定价与补货决策。核心组件——Prophet时序分解、报童随机优化、价格弹性调整——均可直接迁移，仅需根据具体商品的保鲜期调整（单周期$\to$多周期）、需求特征调整（季节性强度）和成本结构调整（损耗率、折扣策略）。

在商业应用层面，该模型可为生鲜零售商提供数据驱动的每日补货与定价建议。Prophet的分解结果可视化有助于管理者理解需求驱动因素，报童模型的优化结果可直接转化为操作指令（各品类/单品的订货量和售价）。随着建议数据（库存记录、折扣原因码、天气数据）的逐步采集，模型预测精度和参数估计可靠性将进一步提升，形成数据采集与模型优化的正向循环。
